\documentclass[a4paper,12pt]{article}
\usepackage[utf8]{inputenc}
\usepackage[T1]{fontenc}
\usepackage[ngerman]{babel}
\usepackage{amsmath}
\usepackage{amssymb}
\usepackage{enumitem}
\usepackage{geometry}
\geometry{a4paper, margin=2.5cm}

\title{Einsendeaufgaben -- Lektion 1\\[0.5em]
\large Modul 61111: Mathematische Grundlagen}
\author{}
\date{}

\begin{document}
\maketitle

\section*{Aufgabe 1.2}

\begin{enumerate}[label=\alph*)]

\item Behauptung: $\displaystyle\sum_{i=1}^{n} i \cdot 2^i = (n-1)\cdot 2^{n+1} + 2$ für alle $n \in \mathbb{N}$.

\textbf{Induktionsanfang} ($n = 1$):
\[
  \sum_{i=1}^{1} i \cdot 2^i = 1 \cdot 2 = 2 = (1-1)\cdot 2^{1+1} + 2 = (n-1)\cdot 2^{n+1} + 2. \checkmark
\]

\textbf{Induktionsschritt} ($n \to n+1$):
\begin{align*}
  \sum_{i=1}^{n+1} i \cdot 2^i
    &= (n+1)\cdot 2^{n+1} + \sum_{i=1}^{n} i \cdot 2^i \\
    &\stackrel{\mathrm{I.V.}}{=} (n+1)\cdot 2^{n+1} + (n-1)\cdot 2^{n+1} + 2 \\
    &= 2^{n+1}\bigl((n+1) + (n-1)\bigr) + 2 \\
    &= 2^{n+1} \cdot 2n + 2 \\
    &= \bigl((n+1)-1\bigr)\cdot 2^{(n+1)+1} + 2. \qed
\end{align*}

\item Behauptung: $6 \mid 7^n - 1$ für alle $n \in \mathbb{N}$.

\textbf{Induktionsanfang} ($n = 1$):
\[
  7^1 - 1 = 6, \quad \text{also gilt } 6 \mid 7^1 - 1. \checkmark
\]

\textbf{Induktionsschritt} ($n \to n+1$): Zu zeigen ist $6 \mid 7^{n+1} - 1$.

Es gilt:
\begin{align*}
  7^{n+1} - 1 &= 7^n \cdot 7 - 1
               = 7^n \cdot 7 - 7 + 6
               = (7^n - 1) \cdot 7 + 6.
\end{align*}
Nach Induktionsvoraussetzung gilt $6 \mid 7^n - 1$.
Da außerdem $6 \mid 6$ gilt, folgt $6 \mid (7^n - 1) \cdot 7 + 6$,
denn allgemein gilt: aus $a \mid b$ und $a \mid c$ folgt $a \mid b + c$
(elementare Zahlentheorie).
Wegen der Gleichheit $(7^n - 1) \cdot 7 + 6 = 7^{n+1} - 1$ ist die Behauptung bewiesen. \qed

\end{enumerate}

\end{document}
